% ============================================================
%  ch03c：§3.5 主结果的表述（局部化；定理 3.5.2.1）（原文 pp.21–22）
%  边界：本块起于 §3.5 标题（p021 中部）。p021 顶部的 3.4.3 与命题 3.4.3.1
%  属 §3.4，归 ch03b。本块止于定理 3.5.2.1（p022 中部）；§3.6 起于 p022
%  下半页，归 ch03d。
% ============================================================
\subsection{主要结果的表述}

\pnum{我们将考察映射 $\Eisspec$ 的\emph{局部化}版本。下面简要讨论相关的形式体系。}

设给定交换图：
\[
\begin{tikzcd}[column sep=huge, row sep=large]
  & \LS^{\arthm}_{\bB} \arrow[dl] \arrow[dr] & \\
  \LS^{\arthm}_{\bG} \arrow[dr, "f"{below}] & & \LS^{\arthm}_{\bT} \arrow[dl, "g"{below}] \\
  & \A^1 &
\end{tikzcd}
\]

令 $\mathring{\A}^1 := \A^1\setminus 0$。注意：
\[
  \Gamma(\LS^{\arthm}_{\bG},\omega)[f^{-1}]
  := \colim_{f\cdot-}\Gamma\bigl(\LS^{\arthm}_{\bG},\omega_{\LS^{\arthm}_{\bG}}\bigr)
  \simeq \Gamma\bigl(\LS^{\arthm}_{\bG}\times_{\A^1}\mathring{\A}^1,\omega\bigr)
\]
类似地，可以用 $(\check T,g)$ 或 $(\check B, f|_{\LS^{\arthm}_{\check B}} = g|_{\LS^{\arthm}_{\check B}})$ 代替 $(\bG,f)$。

现在注意到我们有一个对应：
\begin{equation}
\begin{tikzcd}[column sep=huge, row sep=large]
  & \LS^{\arthm}_{\bB}\times_{\A^1}\mathring{\A}^1 \arrow[dl] \arrow[dr] & \\
  \LS^{\arthm}_{\bG}\times_{\A^1}\mathring{\A}^1 & & \LS^{\arthm}_{\bT}\times_{\A^1}\mathring{\A}^1
\end{tikzcd}
\end{equation}
其中左箭头紧合且可表，右箭头最终余连通且 1-可表。% [Tutor Note: proper 首次括注按 glossary 归 ch03a，此处只用「紧合」。]如同在谱 Eisenstein 级数的定义中那样，我们得到典范映射：
\begin{equation}
  \Gamma\bigl(\LS^{\arthm}_{\bT},\omega_{\LS^{\arthm}_{\bT}}\bigr)[g^{-1}]
  \longrightarrow
  \Gamma\bigl(\LS^{\arthm}_{\bG},\omega_{\LS^{\arthm}_{\bG}}\bigr)[f^{-1}].
\end{equation}
我们显然有：

\begin{lem}
\begin{enumerate}[label=(\arabic*), leftmargin=2.2em]
\item 映射 $\Eisspec:\Gamma(\LS^{\arthm}_{\bT},\omega_{\LS^{\arthm}_{\bT}})\to\Gamma(\LS^{\arthm}_{\bG},\omega_{\LS^{\arthm}_{\bG}})$
把乘以 $g$ 的算子与乘以 $f$ 的算子衔接起来，即该映射自然地是 $\e[t]$-模的态射。
\item 映射 (3.10) 由反转 $t$ 的作用得到，即在 $\e[t]$ 上张量 $\e[t,t^{-1}]$。
\item 图表
\[
\begin{tikzcd}[column sep=huge, row sep=large]
  \Gamma\bigl(\LS^{\arthm}_{\bT},\omega_{\LS^{\arthm}_{\bT}}\bigr)
    \arrow[r, "\Eisspec"] \arrow[d]
  & \Gamma\bigl(\LS^{\arthm}_{\bG},\omega_{\LS^{\arthm}_{\bG}}\bigr) \arrow[d] \\
  \Gamma\bigl(\LS^{\arthm}_{\bT},\omega_{\LS^{\arthm}_{\bT}}\bigr)[g^{-1}]
    \arrow[r, "(3.10)"]
  & \Gamma\bigl(\LS^{\arthm}_{\bG},\omega_{\LS^{\arthm}_{\bG}}\bigr)[f^{-1}]
\end{tikzcd}
\]
交换。
\end{enumerate}
\end{lem}

据此，当含义无歧义时，我们滥用记号，把映射 (3.10) 也记作 $\Eisspec$。

\pnum{现在我们能够陈述关于 $\Eisspec$ 的主要结果。}

\begin{thm}
存在函数 $\delta_{\bG}\in\Om^{\infty}\Gamma(\LS^{\arthm}_{\bG},\OO)$ 与
$\delta_{\bT}\in\Om^{\infty}\Gamma(\LS^{\arthm}_{\bT},\OO)$，它们填入交换图
\begin{equation}
\begin{tikzcd}[column sep=huge, row sep=large]
  & \LS^{\arthm}_{\bB} \arrow[dl] \arrow[dr] & \\
  \LS^{\arthm}_{\bG} \arrow[dr, "\delta_{\bG}"'] & & \LS^{\arthm}_{\bT} \arrow[dl, "\delta_{\bT}"] \\
  & \A^1 &
\end{tikzcd}
\end{equation}
并且满足：
\begin{enumerate}[label=(\arabic*), leftmargin=2.2em]
\item 映射 $\delta_{\bG}$ 在平凡 Weil 局部系统处取非零值。
\item 对象
\[
  \Gamma\bigl(\LS^{\arthm}_{\bG},\omega_{\LS^{\arthm}_{\bG}}\bigr)[\delta_{\bG}^{-1}]\in\Vect
\]
位于上同调次数 $0$，\footnote{注意，(3.8) 预言
$\Gamma(\LS^{\arthm}_{\bG},\omega_{\LS^{\arthm}_{\bG}})$ 位于上同调次数 $0$。尽管我们最终会假设受限几何朗兰兹，
我们在此力求陈述一个不依赖于它的定理，因此保留了这个论断。}以 $\bT$ 代替 $\bG$ 亦类似。
\item 映射
\begin{equation}
  \Eisspec:\Gamma\bigl(\LS^{\arthm}_{\bT},\omega_{\LS^{\arthm}_{\bT}}\bigr)[\delta_{\bT}^{-1}]
  \longrightarrow
  \Gamma\bigl(\LS^{\arthm}_{\bG},\omega_{\LS^{\arthm}_{\bG}}\bigr)[\delta_{\bG}^{-1}]
  \in\Vect^{\heartsuit}
\end{equation}
是满射。
\end{enumerate}
\end{thm}

该结果将在 §4 中证明。
